% ═══════════════════════════════════════════════════════════════
% ARBEITSBLATT: Reaktionsgleichungen aufstellen und ausgleichen
% Thema 2.5 | Chemie Klasse 8 | 36 BE
% ═══════════════════════════════════════════════════════════════

\documentclass[10pt, a4paper]{article}

\usepackage[utf8]{inputenc}
\usepackage[T1]{fontenc}
\usepackage[ngerman]{babel}

\usepackage{helvet}
\renewcommand{\familydefault}{\sfdefault}

\usepackage{microtype}
\usepackage[margin=1.3cm, top=1.8cm, bottom=1.2cm]{geometry}
\usepackage{enumitem}
\usepackage{tabularx}
\usepackage{booktabs}
\usepackage{array}
\usepackage{multicol}
\usepackage{tikz}
\usetikzlibrary{calc, shapes.geometric, arrows.meta}

\usepackage{amsmath, amssymb}
\usepackage{siunitx}
\sisetup{locale = DE, per-mode = symbol}

\usepackage{fancyhdr}
\usepackage{lastpage}

\usepackage{xcolor}
\usepackage{tcolorbox}
\tcbuselibrary{skins}

\usepackage{qrcode}

% ═══════════════════════════════════════════════════════════════
% FARBEN
% ═══════════════════════════════════════════════════════════════

\definecolor{chemie}{HTML}{2a7a4b}
\definecolor{hellgrau}{HTML}{F5F5F5}
\definecolor{mittelgrau}{HTML}{888888}
\definecolor{dunkelgrau}{HTML}{444444}
\definecolor{metall}{HTML}{2c5aa0}
\definecolor{sauerstoff}{HTML}{c62828}

\colorlet{fachfarbe}{chemie}

% ═══════════════════════════════════════════════════════════════
% VARIABLEN
% ═══════════════════════════════════════════════════════════════

\newcommand{\fach}{Chemie}
\newcommand{\thema}{Reaktionsgleichungen aufstellen und ausgleichen}
\newcommand{\klasse}{8}
\newcommand{\maxpunkte}{36}

% ═══════════════════════════════════════════════════════════════
% KOPF- UND FUSSZEILE
% ═══════════════════════════════════════════════════════════════

\pagestyle{fancy}
\fancyhf{}
\renewcommand{\headrulewidth}{0.4pt}
\renewcommand{\footrulewidth}{0pt}
\lhead{\small \fach{} Klasse \klasse{} | \thema}
\rhead{\small Name: \underline{\hspace{3cm}}}
\cfoot{\small\textcolor{mittelgrau}{Seite \thepage{} von \pageref{LastPage}}}

% ═══════════════════════════════════════════════════════════════
% BOXEN
% ═══════════════════════════════════════════════════════════════

\newtcolorbox{merkkasten}{
    colback=hellgrau,
    colframe=hellgrau,
    boxrule=0pt,
    arc=0pt,
    left=6pt, right=6pt, top=5pt, bottom=4pt,
    borderline north={2pt}{0pt}{fachfarbe},
    before skip=4pt, after skip=4pt
}

\newtcolorbox{algorithmusbox}[1][]{
    enhanced,
    colback=fachfarbe!8,
    colframe=fachfarbe,
    boxrule=1.5pt,
    arc=0pt,
    left=8pt, right=8pt, top=12pt, bottom=6pt,
    before skip=10pt, after skip=6pt,
    title={#1},
    fonttitle=\bfseries\small,
    coltitle=white,
    attach boxed title to top left={yshift=-3mm, xshift=4mm},
    boxed title style={colback=fachfarbe, arc=0pt, boxrule=0pt, left=4pt, right=4pt}
}

\newtcolorbox{hinweisbox}{
    colback=metall!10,
    colframe=metall,
    boxrule=1pt,
    arc=0pt,
    left=6pt, right=6pt, top=4pt, bottom=4pt,
    before skip=4pt, after skip=4pt
}

% ═══════════════════════════════════════════════════════════════
% BEFEHLE
% ═══════════════════════════════════════════════════════════════

\newcommand{\luecke}[1]{\underline{\hspace{#1}}}
\newcommand{\punktebox}[1]{\hfill\small\textcolor{mittelgrau}{[\_\_/#1]}}

\newcommand{\aufgabe}[2]{%
    \vspace{6pt}%
    \noindent\textbf{\textcolor{fachfarbe}{Aufgabe #1}} #2%
    \vspace{3pt}%
    \nopagebreak[4]%
}

\newcommand{\operator}[1]{\textbf{\textcolor{fachfarbe}{#1}}}

\newlist{teil}{enumerate}{2}
\setlist[teil,1]{label=\alph*), leftmargin=1.6em, itemsep=2pt, parsep=0pt, topsep=2pt}

\newcommand{\antwortzeilen}[1]{%
    \par\vspace{3pt}%
    \foreach \n in {1,...,#1}{%
        \noindent\rule{\linewidth}{0.3pt}\vspace{5pt}%
    }%
}

\setlength{\parindent}{0pt}
\setlength{\parskip}{2pt}

% ═══════════════════════════════════════════════════════════════
% DOKUMENT
% ═══════════════════════════════════════════════════════════════

\begin{document}

% ═══════════════════════════════════════════════════════════════
% HEADER
% ═══════════════════════════════════════════════════════════════

\begin{center}
    {\large\bfseries\textcolor{fachfarbe}{\thema}}\\[0.1em]
    {\small Arbeitsblatt | \fach{} Klasse \klasse{} | Thema 2.5}
\end{center}
\vspace{-0.4em}
\hrule
\vspace{0.4em}

% ═══════════════════════════════════════════════════════════════
% WERTIGKEIT + TABELLE
% ═══════════════════════════════════════════════════════════════

\begin{minipage}[t]{0.58\textwidth}
\begin{merkkasten}
\textbf{Wertigkeit – Die „Händezahl"}\\[0.2em]
Die \luecke{2.5cm} gibt an, wie viele \luecke{2.5cm} ein Atom eingehen kann.\\[0.3em]
Man kann sie sich als „\luecke{2.5cm}" vorstellen.
\end{merkkasten}
\end{minipage}%
\hfill
\begin{minipage}[t]{0.38\textwidth}
\centering
\small
\textbf{Wertigkeitstabelle}\\[0.3em]
\begin{tabular}{|c|c|c|c|c|c|c|c|}
\hline
\textbf{Na} & \textbf{Mg} & \textbf{Al} & \textbf{Fe} & \textbf{Cu} & \textbf{Zn} & \textbf{Ca} & \textbf{O} \\
\hline
I & II & III & II/III & I/II & II & II & II \\
\hline
\end{tabular}
\end{minipage}

\vspace{4pt}

\begin{hinweisbox}
\small
\textbf{Warum bilden sich diese Formeln?} Bei Metalloxiden gibt das Metall Elektronen ab, Sauerstoff nimmt sie auf. So entstehen \textbf{Ionen}, die sich anziehen (Ionenbindung). Die Wertigkeit zeigt, wie viele Elektronen abgegeben bzw. aufgenommen werden.
\end{hinweisbox}

\vspace{2pt}

% ═══════════════════════════════════════════════════════════════
% ALGORITHMUS 1: VERHÄLTNISFORMEL
% ═══════════════════════════════════════════════════════════════

\begin{algorithmusbox}[Algorithmus 1: Verhältnisformel aufstellen]
\small
\begin{tabular}{@{}c p{12.5cm}@{}}
\textbf{1} & \textbf{Elementsymbole aufschreiben:} \hfill Li \quad O \\[0.4em]
\textbf{2} & \textbf{Wertigkeiten ablesen:} \hfill I \quad II \\[0.4em]
\textbf{3} & \textbf{Kreuzregel anwenden:} Wertigkeit von A → Index von B (und umgekehrt)\\
& \hfill Li\textsubscript{\luecke{0.5cm}}O\textsubscript{\luecke{0.5cm}} → \textbf{Li\textsubscript{2}O} \\
\end{tabular}
\end{algorithmusbox}

\vspace{-2pt}

% ═══════════════════════════════════════════════════════════════
% ALGORITHMUS 2: REAKTIONSGLEICHUNG
% ═══════════════════════════════════════════════════════════════

\begin{algorithmusbox}[Algorithmus 2: Reaktionsgleichung aufstellen und ausgleichen]
\small
\begin{tabular}{@{}c p{12.5cm}@{}}
\textbf{1} & \textbf{Wortgleichung:} Lithium + Sauerstoff → Lithiumoxid \\[0.4em]
\textbf{2} & \textbf{Formelgleichung:} Li + O\textsubscript{2} → Li\textsubscript{2}O \quad {\scriptsize(Achtung: Sauerstoff ist O\textsubscript{2}!)} \\[0.4em]
\textbf{3} & \textbf{Atome zählen:} Links: Li = \luecke{0.5cm}, O = \luecke{0.5cm} \quad Rechts: Li = \luecke{0.5cm}, O = \luecke{0.5cm} \\[0.4em]
\textbf{4} & \textbf{Vorfaktoren ausgleichen} (Atomzahlen unter die Gleichung schreiben): \\[0.5em]
& \textbf{Start:} \begin{tabular}[t]{@{}c@{}c@{}c@{}c@{}c@{}}
Li & $\;+\;$ & O\textsubscript{2} & $\;\rightarrow\;$ & Li\textsubscript{2}O \\[-0.2em]
{\scriptsize\luecke{0.4cm}} & & {\scriptsize\luecke{0.4cm}} & & {\scriptsize\luecke{0.4cm}\hspace{0.6em}\luecke{0.4cm}}
\end{tabular} \quad $\Rightarrow$ nicht gleich! \\[0.8em]
& a) \textbf{Produkt:} \begin{tabular}[t]{@{}c@{}c@{}c@{}c@{}c@{}}
Li & $\;+\;$ & O\textsubscript{2} & $\;\rightarrow\;$ & \luecke{0.4cm}\,Li\textsubscript{2}O \\[-0.2em]
{\scriptsize\luecke{0.4cm}} & & {\scriptsize\luecke{0.4cm}} & & {\scriptsize\luecke{0.4cm}\hspace{0.6em}\luecke{0.4cm}}
\end{tabular} \\[0.8em]
& b) \textbf{Metall:} \begin{tabular}[t]{@{}c@{}c@{}c@{}c@{}c@{}}
\luecke{0.4cm}\,Li & $\;+\;$ & O\textsubscript{2} & $\;\rightarrow\;$ & \luecke{0.4cm}\,Li\textsubscript{2}O \\[-0.2em]
{\scriptsize\luecke{0.4cm}} & & {\scriptsize\luecke{0.4cm}} & & {\scriptsize\luecke{0.4cm}\hspace{0.6em}\luecke{0.4cm}}
\end{tabular} \\[0.8em]
& \textbf{Ergebnis:} \fbox{\luecke{0.4cm} Li + \luecke{0.4cm} O\textsubscript{2} $\rightarrow$ \luecke{0.4cm} Li\textsubscript{2}O} \\
\end{tabular}
\end{algorithmusbox}

% ═══════════════════════════════════════════════════════════════
% AUFGABE 1: WERTIGKEITEN (★)
% ═══════════════════════════════════════════════════════════════

\aufgabe{1 ($\star$)}{Wertigkeiten ablesen} \punktebox{4}

\operator{Lies} die Wertigkeiten aus der Tabelle ab und \operator{ergänze}.

\small
\begin{tabular}{|p{3cm}|p{2cm}|p{8cm}|}
\hline
\textbf{Element} & \textbf{Wertigkeit} & \textbf{Bedeutung} \\
\hline
Magnesium (Mg) & & Mg gibt \luecke{1cm} Elektron(en) ab. \\
\hline
Aluminium (Al) & & Al gibt \luecke{1cm} Elektron(en) ab. \\
\hline
Sauerstoff (O) & & O nimmt \luecke{1cm} Elektron(en) auf. \\
\hline
Calcium (Ca) & & Ca gibt \luecke{1cm} Elektron(en) ab. \\
\hline
\end{tabular}

\vspace{4pt}

% ═══════════════════════════════════════════════════════════════
% AUFGABE 2: VERHÄLTNISFORMELN (★★)
% ═══════════════════════════════════════════════════════════════

\aufgabe{2 ($\star\star$)}{Verhältnisformeln mit der Kreuzregel} \punktebox{6}

\operator{Wende} Algorithmus 1 an und \operator{schreibe} die Oxidformel.

\small
\begin{tabular}{|p{4cm}|p{3cm}|p{5cm}|}
\hline
\textbf{Metall + Sauerstoff} & \textbf{Wertigkeiten} & \textbf{Oxidformel} \\
\hline
Magnesium + Sauerstoff & Mg(II), O(II) & MgO \quad {\scriptsize(gleiche Wertigkeit → 1:1)} \\
\hline
Calcium + Sauerstoff & Ca(\luecke{0.5cm}), O(\luecke{0.5cm}) & \\
\hline
Zink + Sauerstoff & Zn(\luecke{0.5cm}), O(\luecke{0.5cm}) & \\
\hline
Aluminium + Sauerstoff & Al(\luecke{0.5cm}), O(\luecke{0.5cm}) & \\
\hline
Eisen + Sauerstoff & Fe(III), O(II) & \\
\hline
Kupfer + Sauerstoff & Cu(II), O(II) & \\
\hline
\end{tabular}

\newpage

% ═══════════════════════════════════════════════════════════════
% AUFGABE 3: VOLLSTÄNDIGER ALGORITHMUS (★★)
% ═══════════════════════════════════════════════════════════════

\aufgabe{3 ($\star\star$)}{Reaktionsgleichungen Schritt für Schritt} \punktebox{16}

\operator{Wende} den Algorithmus an. \operator{Fülle} die Lücken aus.

\small

\textbf{a) Magnesium – einfacher Fall} (Wertigkeit: II, Oxidformel: MgO)

\begin{tabular}{@{}l@{}}
\textbf{Start:} Atome zählen \\[0.3em]
\hspace{1em}\begin{tabular}[t]{@{}c@{}c@{}c@{}c@{}c@{}}
Mg & $\;+\;$ & O\textsubscript{2} & $\;\rightarrow\;$ & MgO \\[-0.2em]
{\scriptsize\luecke{0.4cm}} & & {\scriptsize\luecke{0.4cm}} & & {\scriptsize\luecke{0.4cm}\hspace{0.5em}\luecke{0.4cm}}
\end{tabular} \quad $\Rightarrow$ nicht gleich! \\[0.8em]
%
a) \textbf{Produkt:} O rechts soll = O links (2) sein \\[0.3em]
\hspace{1em}\begin{tabular}[t]{@{}c@{}c@{}c@{}c@{}c@{}}
Mg & $\;+\;$ & O\textsubscript{2} & $\;\rightarrow\;$ & \luecke{0.4cm}\,MgO \\[-0.2em]
{\scriptsize\luecke{0.4cm}} & & {\scriptsize\luecke{0.4cm}} & & {\scriptsize\luecke{0.4cm}\hspace{0.5em}\luecke{0.4cm}}
\end{tabular} \\[0.8em]
%
b) \textbf{Metall:} Mg links = Mg rechts \\[0.3em]
\hspace{1em}\begin{tabular}[t]{@{}c@{}c@{}c@{}c@{}c@{}}
\luecke{0.4cm}\,Mg & $\;+\;$ & O\textsubscript{2} & $\;\rightarrow\;$ & \luecke{0.4cm}\,MgO \\[-0.2em]
{\scriptsize\luecke{0.4cm}} & & {\scriptsize\luecke{0.4cm}} & & {\scriptsize\luecke{0.4cm}\hspace{0.5em}\luecke{0.4cm}}
\end{tabular} \\[0.6em]
%
\textbf{Ergebnis:} \fbox{\luecke{0.4cm} Mg + \luecke{0.4cm} O\textsubscript{2} $\rightarrow$ \luecke{0.4cm} MgO} \\
\end{tabular}

\vspace{8pt}

\textbf{b) Eisen – komplexer Fall} (Wertigkeit: III, Oxidformel: Fe\textsubscript{2}O\textsubscript{3})

\begin{tabular}{@{}l@{}}
\textbf{Start:} Atome zählen \\[0.3em]
\hspace{1em}\begin{tabular}[t]{@{}c@{}c@{}c@{}c@{}c@{}}
Fe & $\;+\;$ & O\textsubscript{2} & $\;\rightarrow\;$ & Fe\textsubscript{2}O\textsubscript{3} \\[-0.2em]
{\scriptsize\luecke{0.4cm}} & & {\scriptsize\luecke{0.4cm}} & & {\scriptsize\luecke{0.4cm}\hspace{0.5em}\luecke{0.4cm}}
\end{tabular} \quad $\Rightarrow$ nicht gleich! \\[0.8em]
%
a) \textbf{Produkt:} O rechts soll gerade sein (kgV von 2 und 3 = 6) \\[0.3em]
\hspace{1em}\begin{tabular}[t]{@{}c@{}c@{}c@{}c@{}c@{}}
Fe & $\;+\;$ & O\textsubscript{2} & $\;\rightarrow\;$ & \luecke{0.4cm}\,Fe\textsubscript{2}O\textsubscript{3} \\[-0.2em]
{\scriptsize\luecke{0.4cm}} & & {\scriptsize\luecke{0.4cm}} & & {\scriptsize\luecke{0.4cm}\hspace{0.5em}\luecke{0.4cm}}
\end{tabular} \\[0.8em]
%
b) \textbf{O\textsubscript{2}:} O links = O rechts \\[0.3em]
\hspace{1em}\begin{tabular}[t]{@{}c@{}c@{}c@{}c@{}c@{}}
Fe & $\;+\;$ & \luecke{0.4cm}\,O\textsubscript{2} & $\;\rightarrow\;$ & \luecke{0.4cm}\,Fe\textsubscript{2}O\textsubscript{3} \\[-0.2em]
{\scriptsize\luecke{0.4cm}} & & {\scriptsize\luecke{0.4cm}} & & {\scriptsize\luecke{0.4cm}\hspace{0.5em}\luecke{0.4cm}}
\end{tabular} \\[0.8em]
%
c) \textbf{Metall:} Fe links = Fe rechts \\[0.3em]
\hspace{1em}\begin{tabular}[t]{@{}c@{}c@{}c@{}c@{}c@{}}
\luecke{0.4cm}\,Fe & $\;+\;$ & \luecke{0.4cm}\,O\textsubscript{2} & $\;\rightarrow\;$ & \luecke{0.4cm}\,Fe\textsubscript{2}O\textsubscript{3} \\[-0.2em]
{\scriptsize\luecke{0.4cm}} & & {\scriptsize\luecke{0.4cm}} & & {\scriptsize\luecke{0.4cm}\hspace{0.5em}\luecke{0.4cm}}
\end{tabular} \\[0.6em]
%
\textbf{Ergebnis:} \fbox{\luecke{0.4cm} Fe + \luecke{0.4cm} O\textsubscript{2} $\rightarrow$ \luecke{0.4cm} Fe\textsubscript{2}O\textsubscript{3}} \\
\end{tabular}

\vspace{8pt}

\begin{minipage}[t]{0.46\textwidth}
\textbf{c) Aluminiumoxid}

{\small Al hat Wertigkeit III}

\vspace{3.5em}

\textbf{Ergebnis}

\fbox{\luecke{0.4cm} Al + \luecke{0.4cm} O\textsubscript{2} $\rightarrow$ \luecke{0.4cm} \luecke{1.8cm}}
\end{minipage}%
\hfill
\vrule
\hfill
\begin{minipage}[t]{0.46\textwidth}
\textbf{d) Kupferoxid}

{\small Cu hat Wertigkeit II}

\vspace{3.5em}

\textbf{Ergebnis}

\fbox{\luecke{0.4cm} Cu + \luecke{0.4cm} O\textsubscript{2} $\rightarrow$ \luecke{0.4cm} \luecke{1.8cm}}
\end{minipage}

\vspace{8pt}

% ═══════════════════════════════════════════════════════════════
% AUFGABE 4: FEHLERANALYSE (★★★)
% ═══════════════════════════════════════════════════════════════

\aufgabe{4 ($\star\star\star$)}{Fehler finden und korrigieren} \punktebox{6}

Max hat zwei Reaktionsgleichungen aufgestellt. Beide enthalten Fehler.

\operator{Finde} die Fehler, \operator{erkläre} sie und \operator{korrigiere}.

\vspace{4pt}

\textbf{a) Max schreibt:} \quad 2 Al + 3 O → Al\textsubscript{2}O\textsubscript{3}

\vspace{2pt}
Fehler: \antwortzeilen{1}
Korrektur: \luecke{8cm}

\vspace{4pt}

\textbf{b) Max schreibt:} \quad Mg + O\textsubscript{2} → MgO\textsubscript{2}

\vspace{2pt}
Fehler: \antwortzeilen{1}
Korrektur: \luecke{8cm}

\vspace{8pt}

% ═══════════════════════════════════════════════════════════════
% AUFGABE 5: TRANSFER (★★★)
% ═══════════════════════════════════════════════════════════════

\aufgabe{5 ($\star\star\star$)}{Transfer: Unbekanntes Metall} \punktebox{4}

Natrium (Na) hat die Wertigkeit I. \operator{Stelle} die vollständige Reaktionsgleichung auf.

\vspace{4pt}
\begin{tabular}{@{}ll@{}}
Wortgleichung: & Natrium + Sauerstoff → Natriumoxid \\[0.5em]
Oxidformel (Alg. 1): & Na(\luecke{0.5cm}) + O(\luecke{0.5cm}) → Na\textsubscript{\luecke{0.5cm}}O\textsubscript{\luecke{0.5cm}} \\[0.5em]
Ausgeglichene Gleichung (Alg. 2): & \luecke{0.5cm} Na + \luecke{0.5cm} O\textsubscript{2} → \luecke{0.5cm} Na\textsubscript{2}O \\
\end{tabular}

% ═══════════════════════════════════════════════════════════════
% PUNKTE + QR
% ═══════════════════════════════════════════════════════════════

\vfill

\hfill
\begin{minipage}[b]{0.2\textwidth}
\raggedleft
\qrcode[height=1.2cm]{https://devbox42.github.io/sciencesim/chemie/kl08/03-Sauerstoff/02-reaktionsgleichungen/LP-02-reaktionsgleichungen.html}\\
{\tiny Lernpfad}
\end{minipage}

\end{document}
